\begin{abstract}
Weight decay introduces a hyperparameter $\lambda$ that improves stability but disrupts the balance among optimization hyperparameters (learning rate $\eta$, batch size $B$), making optimal tuning substantially more difficult. While prior work has focused on scaling laws for $\eta$ and $B$, $\lambda$ is typically treated as a fixed constant, overlooking its impact on optimization dynamics. In this work, we systematically study how $\lambda$ affects the stability of the optimization system, as well as its coupling with $\eta$ and $B$. Using a stability-based analysis, we derive stability bounds for basic optimizers such as stochastic gradient descent (SGD) and SGD with momentum (SGDM), showing that weight decay can tighten upper bounds but also impose stricter constraints on $\eta$. We further propose a heuristic regularization equivalence hypothesis, under which we derive the analytic coupling of hyperparameters ($\lambda \approx \frac{1}{\eta T}$), consistent with recent empirical results that interpret AdamW as an exponential moving average (EMA). Extending this analysis to the batch size, we show that the optimal product $\lambda\times \eta$ scales with $B$. Experiments across diverse architectures and tasks, from ResNet-18, ResNet-50, and VGG-16 on vision tasks to Qwen3-0.6B on language modeling, validate our theoretical and heuristic findings.

\end{abstract}